% Problema 21 del Cálculo de Spivak
% Demostración de acotación para $|xy - x_0y_0|$

\section*{Problema 21 - Spivak: Demostración de acotación para $|xy - x_0y_0|$}
\addcontentsline{toc}{section}{Problema 21 - Spivak: Demostración de acotación para $|xy - x_0y_0|$}

\subsection*{Enunciado}

Demostrar que si 
\[
|x - x_0| < \min\left(\frac{\varepsilon}{2(|y_0| + 1)}, 1\right) \quad \text{y} \quad |y - y_0| < \frac{\varepsilon}{2(|x_0| + 1)},
\]
entonces $|xy - x_0y_0| < \varepsilon$.

\subsection*{Demostración}

Comenzamos escribiendo $xy - x_0y_0$ de forma estratégica:
\[
xy - x_0y_0 = xy - x_0y + x_0y - x_0y_0 = (x - x_0)y + x_0(y - y_0).
\]

Tomando valor absoluto:
\[
|xy - x_0y_0| = |(x - x_0)y + x_0(y - y_0)| \leq |x - x_0| \cdot |y| + |x_0| \cdot |y - y_0|.
\]

Ahora necesitamos acotar $|y|$. De la primera hipótesis se tiene:
\[
|x - x_0| < 1.
\]

Esto implica:
\[
|x| = |x - x_0 + x_0| \leq |x - x_0| + |x_0| < 1 + |x_0|.
\]

Además, aplicando la desigualdad triangular a $y$:
\[
|y| = |y - y_0 + y_0| \leq |y - y_0| + |y_0|.
\]

Por la segunda hipótesis, $|y - y_0| < \dfrac{\varepsilon}{2(|x_0| + 1)}$, así que:
\[
|y| < \frac{\varepsilon}{2(|x_0| + 1)} + |y_0|.
\]

Ahora acotamos cada término:

\textbf{Primer término:} $|x - x_0| \cdot |y|$

\[
|x - x_0| \cdot |y| < \min\left(\frac{\varepsilon}{2(|y_0| + 1)}, 1\right) \cdot \left(\frac{\varepsilon}{2(|x_0| + 1)} + |y_0|\right).
\]

Por definición del mínimo, $\min \leq \dfrac{\varepsilon}{2(|y_0| + 1)}$:
\[
|x - x_0| \cdot |y| < \frac{\varepsilon}{2(|y_0| + 1)} \cdot \left(\frac{\varepsilon}{2(|x_0| + 1)} + |y_0|\right).
\]

Ahora es clave observar que el factor $\left(\dfrac{\varepsilon}{2(|x_0| + 1)} + |y_0|\right)$ es menor que $(|y_0| + 1)$:
\[
\frac{\varepsilon}{2(|x_0| + 1)} + |y_0| < 1 + |y_0| = |y_0| + 1,
\]
porque $\dfrac{\varepsilon}{2(|x_0| + 1)} < 1$ (ya que $\varepsilon$ es arbitrariamente pequeño y el denominador es $\geq 1$).

Por tanto:
\[
\frac{\varepsilon}{2(|y_0| + 1)} \cdot \left(\frac{\varepsilon}{2(|x_0| + 1)} + |y_0|\right) < \frac{\varepsilon}{2(|y_0| + 1)} \cdot (|y_0| + 1) = \frac{\varepsilon}{2}.
\]

\textbf{Segundo término:} $|x_0| \cdot |y - y_0|$

\[
|x_0| \cdot |y - y_0| < |x_0| \cdot \frac{\varepsilon}{2(|x_0| + 1)} < (|x_0| + 1) \cdot \frac{\varepsilon}{2(|x_0| + 1)} = \frac{\varepsilon}{2}.
\]

(ya que $|x_0| < |x_0| + 1$).

\textbf{Conclusión:}

\[
|xy - x_0y_0| \leq |x - x_0| \cdot |y| + |x_0| \cdot |y - y_0| < \frac{\varepsilon}{2} + \frac{\varepsilon}{2} = \varepsilon.
\]

\subsection*{Observaciones importantes}

\begin{enumerate}
\item La condición $|x - x_0| < 1$ en la primera hipótesis es crucial: garantiza que $x$ no se desvía demasiado de $x_0$, lo que permite acotar $|y|$ en términos de datos iniciales.

\item La división de $\varepsilon$ en dos partes iguales ($\varepsilon/2$ para cada término) es una técnica estándar en demostraciones de límites con dos variables.

\item Este resultado es fundamental para demostrar que el producto de funciones continuas es continuo: si $x \to x_0$ y $y \to y_0$, entonces $xy \to x_0y_0$.

\item La notación $\min(\cdot, \cdot)$ se utiliza para elegir la restricción más fuerte entre dos cotas, asegurando que ambas desigualdades de la hipótesis se satisfagan simultáneamente.
\end{enumerate}
